\documentclass[11pt]{article}
\input{quasse/shared/preamble}
\title{QuaSSE: mathematical and numerical explorations}
\begin{document}
\maketitle

This document holds developing ideas and reasoning worth revisiting. Its first topic is an
experimental representation for trait transport; it is not an implemented replacement for the
production QuaSSE integrator. Established conventions are in the \href{theory.pdf}{theory document}.

\section{Higher spatial order through log-space evolution}\label{sec:log-space}
\subsection{Motivation and status}
One motivation is to improve spatial order without having to prevent negative $D$ or $E$ values
produced by an ordinary-scale discretization. High-order spatial formulas need not preserve
positivity. Evolving a real log value instead represents a positive quantity on exponentiation,
so spatial formulas can be investigated without imposing positivity on their ordinary-scale output.
A second motivation is handling likelihoods spanning a large dynamic range.

This is a representation benefit, not a guarantee of stability or accuracy. For $E$, positivity is
only part of the constraint: $E\leq1$ must also hold. Moreover, exactly zero values are meaningful,
and a finite real logarithm cannot represent them. Existing experiments concern $\log D$ transport;
they do not yet demonstrate a method for evolving the complete coupled $D,E$ system.

\subsection{A worked transformation}
Assume $D(x,t)>0$, once differentiable in time and twice in trait, and suppose for this calculation
that it obeys trait-only transport on an interval where these assumptions hold:
\begin{equation}\label{eq:D-transport}
 D_t=vD_x+\frac q2 D_{xx}.
\end{equation}
Here $v$ and $q$ use the backward convention in the theory document. Setting $h=\log D$ gives
\begin{equation}\label{eq:log-identities}
 D_t=e^h h_t,\qquad D_x=e^h h_x,\qquad
 D_{xx}=e^h(h_{xx}+h_x^2).
\end{equation}
Substitution and division by $e^h>0$ prove the equivalent equation
\begin{equation}\label{eq:log-transport}
 h_t=v h_x+\frac q2\bigl(h_{xx}+h_x^2\bigr).
\end{equation}
The transformation makes a linear transport equation nonlinear. It does not introduce an
approximation at the continuous level; approximations enter when solving it numerically.
A fixed additive offset in $h$ has no effect on this equation. A time-dependent offset would require
an extra time-derivative term and cannot be silently discarded.

For any finite computed real value $h$, $e^h$ is positive in exact arithmetic regardless of whether
the spatial stencil preserves positivity on ordinary-scale inputs. Floating-point exponentiation
can still underflow or overflow. Oscillations, inaccurate values, or failed time integration remain
possible. The same chain-rule identities apply to $k=\log E$ where $E>0$, but the full equation for
$E$ also has reaction terms. In particular, $E(x,0)=0$ when $\rho=1$, and ensuring $k\leq0$ is a
separate issue. A complete method must address these cases before claiming probability preservation.

\subsection{What the existing experiments establish}
The \href{../../validation/log_pde/README.md}{spatial study} compares second-, fourth-, and
sixth-order centered formulas for log transport on Gaussian and non-Gaussian profiles, including
propagation after a node merge. The tested refinements exhibit higher spatial order for the
non-Gaussian cases and can reduce the number of grid points needed for measured accuracy targets.
Gaussian logs alone are inadequate spatial tests: their quadratic form is differentiated exactly
by these stencils in exact arithmetic.

The experiments supply analytic exterior values at every time-integration stage. That isolates
interior spatial behavior while leaving practical boundary treatment unresolved. They omit reaction
and $E$ propagation, and do not establish superiority over the production convolution integrator.
The \href{../../validation/log_pde/temporal/README.md}{temporal study} compares explicit and implicit
solvers at measured error targets. Its findings depend on problem, grid, tolerances, and the extent
of the parameter search; they do not identify one uniformly best solver. Detailed measurements and
reproduction instructions remain in those reports.

\subsection{Alternatives and questions that remain}
Ordinary-scale evolution avoids logarithmic singularities but needs a suitable treatment of
positivity when using higher-order discretizations. Log-space evolution changes that tradeoff:
positivity of the represented value is automatic for finite logs, while nonlinear evolution and
zero values need attention. Existing convolution-based transport is another relevant comparison;
these studies have not made a matched full-likelihood comparison against it.

The remaining questions concern practical boundaries, coupling to diversification and sampling,
handling exactly zero support, preserving $E\leq1$, and performance on full trees. Narrow profiles
also raise time-integration costs even when their spatial representation is simple. These are open
questions, not a prescribed implementation sequence.

\subsection{Development of the question}
The spatial investigation established useful higher-order behavior in controlled transport tests.
A subsequent temporal investigation broadened the solver and tolerance comparison; it showed why
an apparent preference among solvers can depend on search coverage. The retained conclusion is to
consider spatial and temporal costs together. No production replacement has been selected here.
As this question settles, its established mathematical treatment belongs in the theory document;
this exploration can retain a concise account of the choice and meaningful alternatives.
\end{document}
